\chapter{Hindemith, Paul} % (fold)
\label{cha:dohnanyi}

\href{https://en.wikipedia.org/wiki/Paul_Hindemith}{Wikipedia}



\section{Opera}

\noindent{Mörder, Hoffnung der Frauen, in one act, on a libretto by Oskar Kokoschka (1919; premiered 1921)}  \\
\noindent{Das Nusch-Nuschi, in one act, on a libretto by Franz Blei (1920; premiered 1921)} \textit{Queensland SO, Werner Andreas Albert}\marginnote{high}  \\
\noindent{Sancta Susanna, in one act, on a libretto by August Stramm (1921; premiered 1922)}  \\
\noindent{Cardillac, in three acts, on a libretto by Ferdinand Lion after E. T. A. Hoffmann's Das Fräulein von Scuderi (premiered 1926; revised version premiered 1952)}  \\
\noindent{Hin und zurück, an operatic 'sketch' in one scene, on a libretto by Marcellus Schiffer (premiered 1927)}  \\
\noindent{Neues vom Tage, a comic opera (Lustige Oper), on a libretto by Marcellus Schiffer (premiered 1929; revised version premiered 1954)}  \\
\noindent{Lehrstück, music to a Lehrstück in seven scenes by Bertolt Brecht (premiered 1929)} \\
\noindent{Wir bauen eine Stadt, a Schuloper, on a libretto by Robert Seitz (premiered 1930)}  \\
\noindent{Mathis der Maler, in 7 scenes, on a libretto by the composer (1935; premiered 1938)}  \\
\noindent{Die Harmonie der Welt, in 5 acts, on a libretto by the composer (1957; premiered 1957)}  \\
\noindent{The Long Christmas Dinner, in one act, on a libretto by Thornton Wilder (premiered 1961)}  \\


\section{Oratorium}

\noindent{Das Unaufhörliche (1931)}

\noindent{When Lilacs Last in the Dooryard Bloom'd: (A Requiem for those we love), for mezzo-soprano and baritone soloists, chorus, and orchestra, based on the poem by Walt Whitman (1946)}  


\section{Ballet}
\noindent{Der Dämon (1922), Op. 28} \textit{Siegfried Mauser, Frankfurt RSO, Werner Andreas Albert}\marginnote{high} \\
\noindent{Triadisches Ballett, with Oskar Schlemmer (1923)} \\ 
\noindent{Nobilissima Visione, with Léonide Massine (1938)}  \textit{Seattle Symphony, Gerard Schwarz}\marginnote{high} \\
\noindent{The Four Temperaments (scored for piano and strings) (1940)} \textit{Emanuel Ax, Los Angeles Philharmonic, Esa-Pekka Salonen}\marginnote{high} \\
\noindent{Hérodiade, with Martha Graham (1944)} \textit{Siegfried Mauser, Frankfurt RSO, Werner Andreas Albert}\marginnote{high} \\ 
\noindent{Hérodiade, with Martha Graham (1944)} \textit{Annie Gicquel, Frankfurt RSO, Werner Andreas Albert}\marginnote{high} \\ 


\section{Orchestral}
\noindent{Lustige Sinfonietta (1916), Op. 4} \textit{Queensland SO, Werner Andreas Albert}\marginnote{high} \\
\noindent{Ragtime (wohltemperiert)  (1921)}  \textit{BBC PO, Yan Pascal Tortelier}\marginnote{high} \\ 
\noindent{Concerto for Orchestra (1925), Op. 38} \textit{Queensland SO, Werner Andreas Albert}\marginnote{high} \\
\noindent{Konzertmusik for wind orchestra (1926), Op. 41} \\
\noindent{Fünf Stücke für Streichorchester (Five Pieces for String Orchestra) (1927)} \\
\noindent{Konzertmusik for brass and string orchestra (1930), Op. 50} \textit{BBC Scottish SO, Martyn Brabbins}\marginnote{max} \\
\noindent{Philharmonic Concerto (Variations for Orchestra) (1932)} \textit{Queensland SO, Werner Andreas Albert}\marginnote{high} \\
\noindent{Symphony: Mathis der Maler (1933-1934)} \textit{BBC Scottish SO, Martyn Brabbins}\marginnote{max} \\
\noindent{Symphonic Dances (1937)} \textit{BBC PO, Yan Pascal Tortelier}\marginnote{high} \\
\noindent{Nobilissima Visione, suite drawn from the ballet Saint Francis (1938)}  \textit{WDR SO, Marek Janowski}\marginnote{high} \\
\noindent{Symphony in E-flat (1940)} \\
\noindent{Amor und Psyche (Overture to a ballet) (1943)} \textit{Queensland SO, Werner Andreas Albert}\marginnote{high} \\
\noindent{Symphonic Metamorphosis of Themes by Carl Maria von Weber (1943)} \textit{BBC Scottish SO, Martyn Brabbins}\marginnote{max} \\
\noindent{Symphonia Serena (1946)} \textit{Queensland SO, Werner Andreas Albert}\marginnote{high} \\
\noindent{Sinfonietta in E major (1946)} \textit{Queensland SO, Werner Andreas Albert}\marginnote{high} \\
\noindent{Die Harmonie der Welt Symphony (1951)} \textit{Queensland SO, Werner Andreas Albert}\marginnote{high} \\
\noindent{Symphony in B-flat for concert band (1951)} \\
\noindent{Pittsburgh Symphony (1958)}  \textit{BBC PO, Yan Pascal Tortelier}\marginnote{high} \\
\noindent{Marsch über den alten "Schweizerton" (1960)} \textit{Queensland SO, Werner Andreas Albert}\marginnote{high} \\

\section{Concertante}

\noindent{Cello Concerto, Op. 3 (1916)} \\
\noindent{Kammermusik No. 1 for 12 instruments, Op. 24/1 (1921)} \textit{Kronberg Academy Soloists, Schleswig-Holstein Festival Orchestra, Christoph Eschenbach}\marginnote{high} \\
\noindent{Kammermusik No. 2 for piano and 12 instruments, Op. 36/1 (1925)} \textit{Kronberg Academy Soloists, Schleswig-Holstein Festival Orchestra, Christoph Eschenbach}\marginnote{high} \\
\noindent{Kammermusik No. 3 for cello and 10 instruments, Op. 36/2 (1925)} \textit{Kronberg Academy Soloists, Schleswig-Holstein Festival Orchestra, Christoph Eschenbach}\marginnote{high} \\
\noindent{Kammermusik No. 4 for violin and large chamber orchestra, Op. 36/3 (1925)} \textit{Kronberg Academy Soloists, Schleswig-Holstein Festival Orchestra, Christoph Eschenbach}\marginnote{high} \\
\noindent{Kammermusik No. 5 for viola and large chamber orchestra, Op. 36/4 (1927)} \textit{Kronberg Academy Soloists, Schleswig-Holstein Festival Orchestra, Christoph Eschenbach}\marginnote{high} \\
\noindent{Kammermusik No. 6 for viola d'amore and 13 instruments, Op. 46/1 (1927)} \textit{Kronberg Academy Soloists, Schleswig-Holstein Festival Orchestra, Christoph Eschenbach}\marginnote{high} \\
\noindent{Kammermusik No. 7 for organ and 15 instruments, Op. 46/2 (1927)} \textit{Kronberg Academy Soloists, Schleswig-Holstein Festival Orchestra, Christoph Eschenbach}\marginnote{high} \\
\noindent{Klaviermusik mit Orchester for left-hand piano and orchestra, Op. 29 (1923)} \\
\noindent{Konzertmusik for viola and large chamber orchestra, Op. 48 (1930)} \\
\noindent{Konzertmusik for piano, brass and harps, Op. 49 (1930)} \\
\noindent{Konzert für Trautonium in Begleitung des Streichorchesters (1931)} \\
\noindent{Der Schwanendreher for viola and small orchestra (1935)} \\
\noindent{Trauermusik for viola and string orchestra (1936)} \textit{Geraldine Walther, San Francisco Symphnoy, Herbert Blomstedt}\marginnote{high} \\
\noindent{Violin Concerto (1939)} \\
\noindent{Cello Concerto (1940)} \\
\noindent{Piano Concerto (1945)} \textit{Siegfried Mauser, Frankfurt RSO, Werner Andreas Albert}\marginnote{high} \\
\noindent{Clarinet Concerto (1947)}  \textit{Sharon Kam, Frankfurt Radio Symphony, Daniel Cohen}\marginnote{high} \\
\noindent{Horn Concerto (1949)} \\
\noindent{Concerto for Flute, Oboe, Clarinet, Bassoon, Harp and Orchestra (1949)} \\
\noindent{Concerto for Trumpet, Bassoon and Strings (1949)} \\
\noindent{Organ Concerto (1962)} \\



\section{Vocal -- selection}

\noindent{Lustige Lieder in Aargauer Mundart (Merry Songs in the Aargau Dialect), for high voice and piano (1914–16), Op. 5} \\
\noindent{Drei Gesänge, for soprano and large orchestra (1917), Op. 9} \\
\noindent{Melancholie, 4 lieder for mezzo-soprano and string quartet, based on poems by Christian Morgenstern (1919), Op. 13} \\
\noindent{Hymns by Walt Whitman (3), for baritone and piano (1919), Op. 14} \\
\noindent{Acht Gesänge, for soprano voice and piano (1920), Op. 18} \\
\noindent{Des Todes Tod, three songs, based on poems by Eduard Reinacher, for voice, 2 violas and 2 violoncellos (1922), Op. 23a} \\
\noindent{Die junge Magd, six poems by Georg Trakl, for voice, flute, clarinet and string quartet (1922), Op. 23b} \\
\noindent{Tuttifäntchen, Weihnachtsmärchen mit Gesang und Tanz in drei Bildern (Christmas Fairytale with singing and dancing in three scenes)} \\
\noindent{Das Marienleben, song cycle for soprano and piano, based on poems by Rainer Maria Rilke, which exists in two versions (1923/48), Op. 27} \\
\noindent{Sing und Spielmusiken für Liebhaber und Musikfreunde (1928/29), Op. 45} \\
\noindent{"Frau Musika", lyrics by Martin Luther} \\
\noindent{8 canons for voices with instruments} \\
\noindent{"Ein Jäger aus Kurpfalz", for strings and woodwinds} \\
\noindent{"Kleine Klaviermusik", easy pentatonic pieces} \\
\noindent{"Martinslied", soloist or unison choir} \\
\noindent{Hin und zurück, sketch with music, lyrics: Marcellus Schiffer (1927), Op. 45a} \\
\noindent{Six Chansons, 6 pieces for a cappella choir, settings of French poetry by Rainer Maria Rilke (1939)} \\
\noindent{Nine English Songs (1942–43, publ. 1945)} \\
\noindent{Apparebit repentina dies, cantata in four movements for mixed choir (SATB) and brass ensemble (1947)} \\
\noindent{Ite, angeli veloces, cantata (1955)} \\
\noindent{12 Fünfstimmige Madrigale for mixed chorus (1958)} \\
\noindent{Angelus Domini apparuit, motet for soprano or tenor and piano (1958)} \\
\noindent{Mass for mixed chorus (1963)} \\
\noindent{Die Serenaden, for soprano voice, oboe, viola, and cello (1924), Op. 35} \\


\section{Chamber}
\noindent{Kleine Kammermusik for wind quintet (1922), Op. 24/2} \\
\noindent{Ouvertüre zum "Fliegenden Holländer", wie sie eine schlechte Kurkapelle morgens um 7 am Brunnen vom Blatt spielt, for string quartet (c. 1925)} \\
\noindent{Drei Stücke für 5 Instrumente (Three Pieces for Five Instruments), for clarinet in B$\flat$, trumpet in C, violin, contrabass and piano (1925)} \\
\noindent{Duet for viola and cello (1934) (also known as Scherzo)} \\
\noindent{Vier Stücke für Fagott und Violoncello (Four Pieces for Bassoon and Cello) (1941)} \\
\noindent{Eight Pieces for 2 violins, viola, cello and double-bass (1927)} \\
\noindent{Octet for clarinet, bassoon, horn, violin, two violas, cello, and double bass (1958)} \\
\noindent{Quintet for clarinet and string quartet (1923, rev. 1954), Op. 30} \\
\noindent{Quartet for clarinet, violin, cello and piano (1938)} \\
\noindent{String Quartet No. 1 in C (1915), Op. 2} \\
\noindent{String Quartet No. 2 in F minor (1918), Op. 10} \\
\noindent{String Quartet No. 3 in C (1920), Op. 16} \\
\noindent{String Quartet No. 4 (1921), Op. 22} \\
\noindent{String Quartet No. 5 (1923), Op. 32} \\
\noindent{String Quartet No. 6 in E-flat (1943)} \\
\noindent{String Quartet No. 7 in E-flat (1945)} \\
\noindent{String Trio No. 1 (1924), Op. 34} \\
\noindent{String Trio No. 2 (1933)} \\
\noindent{Triosatz [retitled 'Rondo' by editor Siegfried Behrend] for three guitars (1925 or 1930)} \\
\noindent{Trio for viola, heckelphone (or tenor saxophone) and piano (1928), Op. 47} \\
\noindent{Wind Septet for flute, oboe, clarinet, bass-clarinet, bassoon, horn and trumpet (1948)} \\
\noindent{Sonata for Four Horns (1952)} \\
\noindent{Konzertstück for two alto saxophones (1933)} \\
\noindent{"Morgenmusik" Sonata for trumpet, trombone and tuba (1932)} \\
\noindent{Plöner Musiktag (1932): Morgenmusik von Turm zu blasen; Tafelmusik; Mahnung an die Jugend, sich der Musik zu befleissigen (cantata); Abendkonzert (trio for recorders)} \\
\noindent{Des kleinen Elektromusikers Lieblinge, 7 pieces for three trautoniums (1930)} \\


\section{Instrumental}
\noindent{Sonata for Violin and Piano No. 1 in E-flat (1918), Op. 11, No. 1} \\
\noindent{Sonata for Violin and Piano No. 2 in D (1918), Op. 11, No. 2} \\
\noindent{Sonata for Violin and Piano No. 3 in E (1935)} \\
\noindent{Sonata for Violin and Piano No. 4 in C (1939)} \\
\noindent{Sonata for Solo Violin No. 1 (1917), Op. 11, No. 6} \\
\noindent{Sonata for Solo Violin No. 2 (1924), Op. 31, No. 1} \\
\noindent{Sonata for Solo Violin No. 3 (1924), Op. 31, No. 2} \\

\noindent{Sonata for Viola and Piano No. 1 in F major (1919), Op. 11, No. 4} \\
\noindent{Sonata for Viola and Piano No. 2 (1922), Op. 25, No. 4} \\
\noindent{Sonata for Viola and Piano No. 3 in F (1939)} \\
\noindent{Sonata for Solo Viola No. 1 (1919), Op. 11, No. 5} \\
\noindent{Sonata for Solo Viola No. 2 (1922), Op. 25, No. 1} \\
\noindent{Sonata for Solo Viola No. 3 (1923), Op. 31, No. 4} \\
\noindent{Sonata for Solo Viola No. 4 (1937)} \\

\noindent{Kleine Sonate for Viola d'amore and Piano (1922), Op. 25, No. 2} \\

\noindent{Variations on "A Frog He Went a Courting" for Cello and Piano (1941)} \\
\noindent{Three Pieces for Cello and Piano (1917), Op. 8} \\
\noindent{Sonata for Cello and Piano No. 1 (1919), Op. 11, No. 3} \\
\noindent{Sonata for Cello and Piano No. 2 (1948)} \\
\noindent{Kleine Sonate for Cello and Piano (1942)} \\
\noindent{Solo Sonata for Cello (1923), Op. 25, No. 3} \\
\noindent{Drei leichte Stücke for Cello and Piano (1938)} \\
\noindent{Sonata for Double Bass and Piano (1949)} \\
\noindent{Pieces for unaccompanied Double Bass (1929)} \\

\noindent{Canonic Sonatina for Two Flutes (1923), Op. 31, No. 3} \\
\noindent{Eight Pieces for Solo Flute (1927)} \\
\noindent{Sonata for Flute and Piano (1936)} \\
\noindent{Sonata for Oboe and Piano (1938)} \\
\noindent{Sonata for English Horn and Piano (1941)} \\
\noindent{Sonata for Clarinet and Piano (1939)} \\
\noindent{Sonata for Bassoon and Piano (1938)} \\

\noindent{Sonata for French Horn and Piano (1939)} \\
\noindent{Sonata for Four Horns (1952)} \\
\noindent{Alto Horn Sonata (also for Horn or Alto Saxophone) (1943)} \\
\noindent{Sonata for Trumpet and Piano (1939)} \\
\noindent{Sonata for Trombone and Piano (1941)} \\
\noindent{Sonata for Tuba and Piano (1955)} \\
\noindent{Sonata for Harp (1939)} \\


\section{Piano}

\noindent{Eight Waltzes ("Drei wunderschöne Mädchen im Schwarzwald") for piano 4 hands (1916), Op. 6} \\
\noindent{In einer Nacht... (1917–1919), Op. 15} \\
\noindent{Sonate (1920), Op. 17} \\
\noindent{Tanzstücke (1920), Op. 19} \\
\noindent{Suite "1922" (1922), Op. 26} \\
\noindent{Übung in drei Stücken, Part 1 (1924–25), Op. 37} \\
\noindent{Reihe Kleine Stücke, Part 2 (1926), Op. 37} \\
\noindent{Kleine Klaviermusik (1929)} \\
\noindent{Wir bauen eine Stadt (1931)} \\
\noindent{Piano Sonata No. 1 (1936)} \\
\noindent{Piano Sonata No. 2 (1936)} \\
\noindent{Piano Sonata No. 3 (1936)} \\
\noindent{Variations (1936)} \\
\noindent{Piano Sonata (four hands) (1938)} \\
\noindent{Sonata for Two Pianos (four hands) (1942)} \\
\noindent{Ludus Tonalis (1942)} \\


\section{Organ}

\noindent{Two Pieces for Organ (1918)} \\
\noindent{Organ Sonata No. 1 (1937)} \\
\noindent{Organ Sonata No. 2 (1937)} \\
\noindent{Organ Sonata No. 3 (on ancient folk songs) (1940)} \\

